\documentclass[11pt,letterpaper]{article}
\usepackage[margin=0.72in]{geometry}
\usepackage[T1]{fontenc}
\usepackage[utf8]{inputenc}
\usepackage{lmodern}
\usepackage{microtype}
\usepackage{xcolor}
\usepackage{enumitem}
\usepackage{tabularx}
\usepackage{hyperref}

\definecolor{accent}{HTML}{00419B}
\definecolor{muted}{HTML}{555555}

\hypersetup{
  colorlinks=true,
  urlcolor=accent,
  linkcolor=accent
}

\setlength{\parindent}{0pt}
\setlength{\parskip}{0.45em}
\setlist[itemize]{leftmargin=1.25em, topsep=0.2em, itemsep=0.15em}

\newcommand{\MauriName}{Maurizio Ungaro}
\newcommand{\MauriTitle}{Staff Scientist, Jefferson Lab}
\newcommand{\MauriEmail}{\href{mailto:ungaro@jlab.org}{ungaro@jlab.org}}
\newcommand{\MauriHomepage}{\href{https://maureeungaro.github.io/home/}{maureeungaro.github.io/home}}
\newcommand{\MauriGitHub}{\href{https://github.com/maureeungaro}{github.com/maureeungaro}}
\newcommand{\MauriScholar}{\href{https://scholar.google.com/citations?user=zkWYILYAAAAJ\&hl=en}{Google Scholar}}
\newcommand{\MauriInspire}{\href{https://inspirehep.net/authors/1322331}{INSPIRE}}
\newcommand{\MauriOrcid}{\href{https://orcid.org/0000-0001-6982-3310}{ORCID 0000-0001-6982-3310}}
\newcommand{\MauriLinkedIn}{\href{https://www.linkedin.com/in/maurizio-ungaro-37992062/}{LinkedIn}}
\newcommand{\MauriJLab}{\href{https://www.jlab.org}{Jefferson Lab}}
\newcommand{\MauriHallB}{\href{https://www.jlab.org/physics/hall-b}{Jefferson Lab Hall-B}}
\newcommand{\MauriRPI}{\href{https://www.rpi.edu}{Rensselaer Polytechnic Institute}}
\newcommand{\MauriUniGe}{\href{https://unige.it/}{Universita degli Studi di Genova}}
\newcommand{\MauriGEMC}{\href{https://gemc.github.io/home/}{GEMC}}
\newcommand{\MauriGeantFourJLab}{\href{https://jeffersonlab.github.io/g4home/}{Geant4 at JLab}}
\newcommand{\MauriClasTwelve}{\href{https://github.com/gemc/clas12Tags}{CLAS12 Simulations}}
\newcommand{\MauriOSG}{\href{https://maureeungaro.github.io/home/osg/osg}{CLAS12 on OSG}}
\newcommand{\MauriLTCC}{\href{https://maureeungaro.github.io/home/ltcc/ltcc}{Low Threshold Cherenkov Counter}}
\newcommand{\MauriSPringEight}{\href{https://www.spring8.or.jp/en/}{SPring-8}}
\newcommand{\MauriHPS}{\href{https://www.jlab.org/physics/hall-b/other}{Heavy Photon Search}}

\newcommand{\DocHeader}[1]{
  {\LARGE\bfseries \MauriName}\hfill {\color{muted}#1}\\
  Staff Scientist, \MauriJLab{}\\
  \MauriEmail{} \quad
  \MauriHomepage{} \quad
  \MauriGitHub{}\\
  \MauriScholar{} \quad
  \MauriInspire{} \quad
  \MauriOrcid{}\\
  \vspace{0.4em}
  \hrule
  \vspace{0.8em}
}

\newcommand{\RoleEntry}[4]{
  \textbf{#1}\hfill {\color{muted}#2}\\
  #3\\
  #4
}

\newcommand{\ProjectEntry}[3]{
  \textbf{#1}: #2\\
  {\small\href{#3}{#3}}
}


\begin{document}
\DocHeader{Biography}

\section*{Short Bio}
Maurizio Ungaro is a nuclear physicist and simulation software developer in
\MauriHallB{}. His work spans \MauriGEMC{} and GEMC3, Geant4-based detector
simulation, CLAS12 systems, and nucleon-structure research.

\section*{Medium Bio}
Maurizio Ungaro is a Staff Scientist in \MauriHallB{}, working at the
intersection of nuclear physics, detector systems, and scientific software.
He develops \MauriGEMC{} and GEMC3, Python tools through pygemc, and CLAS12
simulation workflows. His recent work includes CAD geometry repair and
validation, reproducible performance studies, and interactive visualization.
He supports distributed simulation production on the Open Science Grid and
works on the CLAS12 Low Threshold Cherenkov Counter. He also writes tutorials
and technical documentation for researchers and students.

\section*{Long Bio}
Maurizio Ungaro is a nuclear physicist and simulation software developer at
\MauriJLab{}. His research focuses on the internal structure and dynamics of
the nucleon, including meson electroproduction, nucleon resonances, and the transition between
hadronic and partonic degrees of freedom. His technical work connects detector
operation, Geant4 simulation, CLAS12 production workflows, and scientific
software infrastructure.

He develops and supports \MauriGEMC{}, a database-driven Geant4 simulation
framework, and its next-generation GEMC3 framework with Python tools through
pygemc. Recent development includes CLAS12 geometry and detector-plugin
migration, CAD mesh repair and optimization, and reproducible construction,
overlap, and particle-tracking checks. He also works on interactive detector
visualization, including Geant4 XR integration, multilingual software tools,
and portable build and installation workflows.

His production work includes CLAS12 simulation releases and simGrid workflows
using HTCondor and Open Science Grid resources. His detector responsibilities
include operation, maintenance, calibration, and performance studies of the
CLAS12 \MauriLTCC{}. Through \MauriGeantFourJLab{}, tutorials, examples, and
technical notes, he helps researchers and students use and validate detector
simulations.

His earlier and collaborative work includes experiments at \MauriSPringEight{}
and participation in the \MauriHPS{} collaboration.

\section*{ORCID / LinkedIn Draft}
\textbf{ORCID biography}: Staff Scientist at \MauriJLab{} working on nuclear
physics, Geant4 simulations, \MauriGEMC{} and GEMC3, pygemc, CLAS12 detector
systems, validated CAD geometry, and distributed scientific computing workflows.

\textbf{LinkedIn headline}: Nuclear Physicist | Geant4 and GEMC Simulation
Developer | CLAS12 Software and Detector Systems | Jefferson Lab

\end{document}
